% Generated from locale/tr/content/history/set-theory/pathologies.tex; do not hand-edit.


\olfileid[tr]{his}{set}{pathology}
\olsection{Patolojiler}

Ancak \emph{limit} tanımının bazı oldukça ``patolojik'' kuruluşlara izin
verdiği ortaya çıktı.

Bolzano 1830'lar dolayında \emph{her yerde sürekli}, fakat \emph{hiçbir yerde
türevlenebilir olmayan} bir fonksiyon keşfetti. (Ne yazık ki Bolzano bunu hiç
yayımlamadı; matematikçiler bu fikirle ilk kez, Weierstrass'ın aynı fikri
bağımsız olarak keşfetmesi sayesinde 1872'de karşılaştılar.)\footnote{Bu
tarihçe, Wikipedia'nın
\href{http://en.wikipedia.org/wiki/Weierstrass_function}{Weierstrass
fonksiyonu} maddesindeki son derece ayrıntılı dipnotlarda belgelenmiştir.}
Bu, en hafif deyişle oldukça şaşırtıcıydı. $|x|$ gibi her yerde sürekli,
fakat belirli bir noktada türevlenebilir olmayan fonksiyonlar bulmak kolaydır.
Ama her yerde sürekli olduğu hâlde \emph{hiçbir yerde} türevlenebilir olmayan
bir fonksiyon bambaşka bir yaratıktır. Bir an için böyle bir fonksiyonu nasıl
çizmeye çalışacağınızı düşünün. Sürekli olmasını sağlamak için onu kaleminizi
kâğıttan hiç kaldırmadan çizebilmeniz gerekir; fakat hiçbir yerde
türevlenebilir olmamasını sağlamak için kaleminizin yönünü durmaksızın ve
ansızın değiştirmeniz gerekir.

Bunu başka ``patolojiler'' izledi. Cantor, 5 Ocak 1874'te Dedekind'e yazdığı
bir mektupta şu problemi ortaya attı:
\begin{quote}
Bir yüzey (örneğin sınırıyla birlikte bir kare) ile bir doğru (örneğin uç
noktalarıyla birlikte bir doğru parçası), yüzeyin her noktasına doğrunun bir
noktası ve tersine doğrunun her noktasına yüzeyin bir noktası karşılık
gelecek biçimde bire bir eşlenebilir mi?

Şu anda bu sorunun yanıtı bana hâlâ çok güç görünüyor---gerçi burada da
insanın \emph{hayır} demeye öylesine güçlü bir eğilimi var ki ispatı neredeyse
gereksiz saymak istiyor. [\citealt{Gouvea2011}'den alıntı]
\end{quote}
Fakat Cantor 1877'de yanıldığını ispatladı. Gerçekte bir doğru ile bir kare
tam olarak aynı sayıda noktaya sahiptir. Dedekind'e 29 Haziran 1877'de
``\emph{je le vois, mais je ne le crois pas}''; yani ``Görüyorum ama
inanmıyorum'' diye yazdı. Matematiğin ``yerleşik tarihinde'' bu söz, yeni
sonuçların dönemin matematikçilerine ne denli \emph{kelimenin tam anlamıyla
inanılmaz} geldiğinin göstergesi olarak sık sık ele alınır. (Yazışma
\citet{Gouvea2011}'de sunulur ve ona \olref[his][set][mythology]{sec}'te
yeniden döneceğiz. Cantor'un ispatı \olref[his][set][cantorplane]{sec}'te
özetlenir.)

Cantor'un sonucundan esinlenen Peano, bir doğrunun bir düzlem üzerine
\emph{pürüzsüzce} eşlenmesinin mümkün olup olmadığını düşünmeye başladı. Bu,
\emph{uzayı dolduran bir eğri} olacaktı. Peano \citeyear{Peano1890}'da tam da
böyle bir eğri kurdu. Bu gerçekten sezgiye aykırıdır: Öklid doğruyu ``genişliği
olmayan uzunluk'' (I. Kitap, 2. Tanım) diye tanımlamıştı; fakat Peano bir
doğruyu uygun biçimde kıvırarak uzunluğunun genişliğe dönüştürülebileceğini
göstermişti. Hilbert \citeyear{Hilbert1891}'da, kendisini açıklayan bazı
resimlerle birlikte, biraz daha sezgisel bir uzay dolduran eğri betimledi.
Eğri aşamalar hâlinde kurulur; kuruluşun ilk altı aşaması şöyledir:
\begin{center}
\begin{tikzpicture}
\begin{scope}[xshift=20pt, yshift=20pt]
	\draw [oldiagcolorC, l-system={Hilbert curve, step=40pt, angle=90, axiom=L, order=1}] lindenmayer system;
\end{scope}
\begin{scope}[xshift=110pt, yshift=10pt]
	\draw [oldiagcolorC, l-system={Hilbert curve, step=20pt, angle=90, axiom=L, order=2}] lindenmayer system;
\end{scope}
\begin{scope}[xshift=205pt, yshift=5pt]
	\draw [oldiagcolorC, l-system={Hilbert curve, step=10pt, angle=90, axiom=L, order=3}] lindenmayer system;
\end{scope}
\begin{scope}[xshift=2.5pt, yshift=-97.5pt]
	\draw [oldiagcolorC, l-system={Hilbert curve, step=5pt, angle=90, axiom=L, order=4}] lindenmayer system;
\end{scope}
\begin{scope}[xshift=101.25pt, yshift=-98.75pt]
	\draw [oldiagcolorC, l-system={Hilbert curve, step=2.5pt, angle=90, axiom=L, order=5}] lindenmayer system;
\end{scope}
\begin{scope}[xshift=200.625pt, yshift=-99.375pt]
	\draw [oldiagcolorC, l-system={Hilbert curve, step=1.25pt, angle=90, axiom=L, order=6}] lindenmayer system;
\end{scope}
\draw[gray] (0pt,0pt) rectangle (80pt, 80pt);
\draw[gray] (100pt,0pt) rectangle (180pt, 80pt);
\draw[gray] (200pt,0pt) rectangle (280pt, 80pt);
\draw[gray] (0pt,-100pt) rectangle (80pt, -20pt);
\draw[gray] (100pt,-100pt) rectangle (180pt, -20pt);
\draw[gray] (200pt,-100pt) rectangle (280pt, -20pt);
\end{tikzpicture}  
\end{center}
O zamana dek kesin bir tanıma kavuşmuş olan limitte karenin tamamı düz
!!{colorC} renkle dolar. Ayrıca Hilbert eğrisi her yerde sürekli, fakat hiçbir
yerde türevlenebilir değildir; sezgisel olarak bunun nedeni sonsuz limitte
fonksiyonun her an birdenbire yön değiştirmesidir. (Hilbert'in kuruluşunu
\olref[his][set][hilbertcurve]{sec}'te daha ayrıntılı olarak özetleyeceğiz.)

İyi ya da kötü, bu ``patolojik'' geometrik kuruluşlar geometrik sezgiye
başvurudan kuşkulanmak için bir neden sayıldı. Küme teorisiyle doruğa ulaşacak
yeni bir matematik yapma tarzının neredeyse \emph{propagandasına} dönüştüler.
Alanın daha sonraki mit yaratımında bu sonuçların hem tümüyle kesin hem de
tümüyle sarsıcı olduğu sık sık yinelendi. Böylece ikili bir amaca hizmet
ettiler: geometrik sezgiye güvenmeye karşı bir uyarı ve yeni düşünme
biçimlerinin verimliliğinin bir gösterimi.

